\documentclass{article}

\usepackage{amsmath}
\usepackage{amssymb}

\newcommand{\ZZ}{\mathbb{Z}}
\newcommand{\RR}{\mathbb{R}}
\newcommand{\CC}{\mathbb{C}}

\begin{document}

{\bf Worksheet \#10; date: 010/01/2018}

{\bf MATH 55 Discrete Mathematics}

\begin{enumerate}
\item Using Fermat's little theorem to find the inverse of $5$ modulo $17$.

\item {\em (Rosen 4.4.39)}
\begin{enumerate}
\item Use Fermat's little theorem to compute $5^{2003} \mod 7$, $5^{2003} \mod 11$ and $5^{2003} \mod 13$.
\item Use your results from part (a) and the Chinese remainder theorem to find $5^{2003} \mod 1001$. (Note that $1001 = 7 \cdot 11 \cdot 13$.)
\end{enumerate}

\item {\em (Rosen 4.4.23)} Show that we can easily factor $n$ when we know that $n$ is the product of two primes, $p$ and $q$, and we know the value of $(p-1)(q-1)$.

\item {\em (Rosen 4.4.25; modified)} Encrypt the message UPLOAD using the RSA system with $n = 53 \cdot 61$ and $e = 17$, translating each letter into integers and grouping together pairs of integers, as done in Example 8.

\item Knowing that $n = 53 \cdot 61$. What is the decryption key in the question above?
\end{enumerate}

\end{document}